% ===================================================================
%  CUADRO N.º 5 — La Casa y su construcción.
%  Traduction 2026 d'apres texto/io, controlee sur texto/fr et
%  texto/en. Consigne : voir texto/es/10-cuadro-01.tex.
%
%  Le tableau 5 est un DIALOGUE : on garde \textsc{} pour le nom du
%  parleur, que le rendu laisse tomber en gardant le texte. Les
%  renvois du plan sont des LETTRES, (a) a (m), et l'on ecrit ce que
%  l'ido ecrit, « (1) » compris la ou l'imprimeur a compose un 1 pour
%  un l : la page de lecture a deja de quoi trancher.
% ===================================================================

%%K t05-apar-1 apar
\VUsaut{6.0mm}
\VUtitre{15.0pt}{60}{CUADRO N.º 5}
\VUsaut{1.4mm}
\VUtitre{11.4pt}{0}{La Casa y su construcción.}
\VUsaut{2.2mm}

%%K t05-tit-1 sub
\VUpk{10.2pt}{40}{Un buen encuentro.}
\VUsaut{3.0mm}

%%K t05-01-1 p
\textsc{Juan}. --- Tío, me gustaría mucho saber cómo se construye una
\VUgras{casa}.

%%K t05-01-2 p
\textsc{El Tío} \textsuperscript{(80)}. --- Es muy fácil, hijo mío; ven conmigo y
te enseñaré la que el señor Bernardo \textsuperscript{(79)} está haciendo
construir y en la que voy a vivir.

%%K t05-01-3 p
\textsc{Juan}. --- ¿Y quién ha dibujado el \VUgras{plano} \textsuperscript{(76)} de
esta casa?

%%K t05-01-4 p
\textsc{El Tío}. --- El \VUgras{arquitecto} \textsuperscript{(75)}; presentó un dibujo
muy claro, que fue aprobado, y el \VUgras{contratista} \textsuperscript{(77)}
empezó las obras.

%%K t05-01-5 p
\textsc{Juan}. --- ¿Y quién es el contratista?

%%K t05-01-6 p
\textsc{El Tío}. --- Es el señor Blanco, el padre de tu amigo Jacobo.
Dirige a los obreros con el señor Rojo, el arquitecto.

%%K t05-01-7 p
¡Vaya! Aquí llega justamente el señor Blanco con su \VUgras{regla}
\textsuperscript{(78)}, y el señor Rojo con su plano.

%%K t05-01-8 p
Buenos días, señores. ¿Cómo están ustedes?

%%K t05-01-9 p
\textsc{El Sr. Rojo}. --- Muy bien, gracias. Buenos días, Juan; ¡cómo ha
crecido este niño!

%%K t05-01-10 p
\textsc{El Tío}. --- Ya lo creo, es todo un hombrecito. Es muy serio; hay
que explicarle todo lo que ve. Hoy querría saber cómo se construye una
casa.

%%K t05-tit-2 sub
\VUpk{10.2pt}{40}{Los obreros.}
\VUsaut{3.0mm}

%%K t05-01-11 p
\textsc{El Sr. Blanco}. --- Entonces ven conmigo, amiguito, y te lo
explicaré enseguida, porque me gustan mucho los niños que quieren
aprender.

%%K t05-01-12 p
Ves a ese obrero que tiene una \VUgras{pala} \textsuperscript{(82)} en la mano: es
un \VUgras{peón caminero} \textsuperscript{(81)}. Está abriendo una \VUgras{zanja}
\textsuperscript{(46)} para colocar la cañería del agua. También fue él quien
cavó el suelo en el sitio donde está la casa, para que el albañil
levantara las cuatro \VUgras{paredes}. La parte de esas paredes que queda
bajo tierra se llama los \VUgras{cimientos} \textsuperscript{(2)}. El espacio
comprendido entre los cimientos forma la \VUgras{bodega} \textsuperscript{(3)}. Esa
bodega tiene una ventanilla llamada \VUgras{tragaluz} \textsuperscript{(5)}, una
\VUgras{puerta} \textsuperscript{(4)} y una escalera.

%%K t05-01-13 p
El \VUgras{albañil} \textsuperscript{(108)} siguió levantando las paredes dejando
huecos para las \VUgras{puertas} \textsuperscript{(8)} y las \VUgras{ventanas}
\textsuperscript{(17)}.

%%K t05-01-14 p
\textsc{Juan}. --- ¡Pero no lo veo, a ese albañil!

%%K t05-01-15 p
\textsc{El Sr. Blanco}. --- Ya ha terminado de trabajar en la casa. Está
junto a la \VUgras{cuadra} \textsuperscript{(54)}, sobre su \VUgras{andamio}
\textsuperscript{(110)}, levantando la pared del \VUgras{lavadero} \textsuperscript{(56)}.

%%K t05-01-16 p
Tiene una llana, una artesa y una \VUgras{plomada} \textsuperscript{(109)}. Pero no
te acordarás de esos nombres.

%%K t05-01-17 p
\textsc{Juan}. --- Sí que me acordaré, porque voy a pedirle a mi tío una
llana pequeña y una artesa, y yo mismo haré una plomada con un cordel.

%%K t05-tit-3 sub
\VUsaut{3.0mm}
\VUpk{10.2pt}{40}{Los materiales.}
\VUsaut{3.0mm}

%%K t05-01-18 p
\textsc{El Tío}. --- ¡Ah, muy bien! Pero dime, Juan, ¿qué hace falta para
construir una casa?

%%K t05-01-19 p
\textsc{Juan}. --- Hacen falta \VUgras{sillares} \textsuperscript{(59)} y
\VUgras{mortero} \textsuperscript{(65)}.

%%K t05-01-20 p
\textsc{El Tío}. --- Exacto. Mira a ese hombre del sombrero grande, en su
\VUgras{carro} \textsuperscript{(136)}. Es un \VUgras{carretero} \textsuperscript{(135)}.
Todos los días trae aquí \VUgras{bloques de piedra} \textsuperscript{(60)}. Descarga
su carro en el patio, y los \VUgras{canteros} \textsuperscript{(83)} labran los
bloques y los sierran con su \VUgras{sierra de piedra} \textsuperscript{(85)}. Un
\VUgras{peón} \textsuperscript{(146)} se los lleva al albañil, que los coloca.

%%K t05-01-21 p
Mientras tanto, el \VUgras{amasador} \textsuperscript{(87)} prepara el mortero.
Necesita \VUgras{arena} \textsuperscript{(62)}, \VUgras{cal} \textsuperscript{(63)} y
\VUgras{agua} \textsuperscript{(64)}. La cal está a su lado en un \VUgras{saco}
\textsuperscript{(71)}; un \VUgras{ayudante de albañil} \textsuperscript{(111)} le trae la
arena en una \VUgras{carretilla} \textsuperscript{(112)}, y el \VUgras{aprendiz}
\textsuperscript{(113)} está en la bomba \textsuperscript{(58)}. Llena un \VUgras{cubo}
\textsuperscript{(114)}. El mortero estará listo enseguida. El
\VUgras{peón de cuezo} \textsuperscript{(107)} lo coge y lo sube por la escalera
hasta el andamio, donde un albañil termina ese lado de la casa.

%%K t05-01-22 p
Y ahora, ¿cómo se llama el obrero que hace el \VUgras{tejado}
\textsuperscript{(31)}?

%%K t05-01-23 p
\textsc{Juan}. --- Es el \VUgras{tejador} \textsuperscript{(118)}.

%%K t05-tit-4 sub
\VUsaut{3.0mm}
\VUpk{10.2pt}{40}{El tejado de la casa.}
\VUsaut{3.0mm}

%%K t05-01-24 p
\textsc{El Sr. Blanco}. --- Sí, pero primero viene el \VUgras{carpintero}
\textsuperscript{(115)}.

%%K t05-01-25 p
El carpintero hace la \VUgras{armadura} \textsuperscript{(35)}. Une las
\VUgras{vigas} \textsuperscript{(37)} con \VUgras{clavijas} \textsuperscript{(117)}. Esos
obreros de allí arriba, con sus anchos calzones de terciopelo, son
carpinteros. Uno sostiene una viga pequeña y el otro corta una clavija
con su \VUgras{hacha} \textsuperscript{(116)}. El que está junto a la
\VUgras{chimenea} \textsuperscript{(41)} es el tejador. Clava listones sobre las
\VUgras{viguetas} (*), y \VUgras{pizarras} \textsuperscript{(66)} sobre los listones.
A veces pone tejas.

%%K t05-noto-1 noto
\VUnotes{143.2mm}{%
(*) \VUgras{Balk-o} = alemán \textit{Balken}, inglés \textit{joist},
francés \textit{solive}, italiano \textit{travicello}, ruso
\textit{balka}, español y portugués \textit{viga}. Raíz técnica todavía
no adoptada, pero que se propone aquí para traducir el sentido de la
palabra francesa \textit{solive}, que no es un \textit{trabo} (francés
\textit{poutre}) ni una viga pequeña. Es un madero escuadrado que
sostiene un suelo y un techo.%
}

%%K t05-01-26 p
El tejado de la cuadra es de \VUgras{tejas} \textsuperscript{(55)}, el de la casa
es de \VUgras{pizarra} \textsuperscript{(32)} y el de la galería es de
\VUgras{cinc} \textsuperscript{(24)}.

%%K t05-01-27 p
\textsc{Juan}. --- ¿El tejador hace también los tejados de cinc?

%%K t05-01-28 p
\textsc{El Sr. Blanco}. --- No, ese es el \VUgras{cinquero} \textsuperscript{(121)} o
el \VUgras{fontanero}, el que está colocando una \VUgras{buhardilla}
\textsuperscript{(38)} en el tejado de la casa. También pone los \VUgras{canalones}
\textsuperscript{(43)} y las \VUgras{bajantes} \textsuperscript{(42)}. Suelda el cinc y la
hojalata. Fue él quien fijó el \VUgras{pararrayos} \textsuperscript{(40)} y la
\VUgras{veleta} \textsuperscript{(39)} en el \VUgras{hastial} \textsuperscript{(36)}.

%%K t05-01-29 p
\textsc{Juan}. --- ¿Entonces la casa está terminada?

%%K t05-tit-5 sub
\VUsaut{3.0mm}
\VUpk{10.2pt}{40}{Las herramientas.}
\VUsaut{3.0mm}

%%K t05-01-30 p
\textsc{El Sr. Blanco}. --- No del todo. El \VUgras{yesero} \textsuperscript{(99)}
levanta con \VUgras{ladrillos} \textsuperscript{(61)} los tabiques que la dividen en
las distintas habitaciones. Reviste de \VUgras{yeso} \textsuperscript{(70)} los
\VUgras{techos} \textsuperscript{(26)} y las paredes. Ahora está haciendo el techo
de la \VUgras{galería} \textsuperscript{(33)}. Tiene, como el albañil, una
\VUgras{llana} \textsuperscript{(100)} y una \VUgras{artesa} \textsuperscript{(101)}.

%%K t05-01-31 p
Después el ebanista hace las puertas, las ventanas, los
\VUgras{entarimados} \textsuperscript{(16)} y las \VUgras{contraventanas}
\textsuperscript{(19)}.

%%K t05-01-32 p
Ahora trabaja en el patio en su \VUgras{banco de carpintero} \textsuperscript{(90)}.
Tiene una \VUgras{sierra} \textsuperscript{(94)} para serrar la \VUgras{madera}
\textsuperscript{(69)}, un \VUgras{cepillo} \textsuperscript{(91)}, un \VUgras{sargento}
\textsuperscript{(92)}, un \VUgras{formón} \textsuperscript{(94 \textit{bis})}, un
\VUgras{compás} \textsuperscript{(95)} y un \VUgras{mazo de madera}
\textsuperscript{(95 \textit{bis})}.

%%K t05-01-33 p
\textsc{Juan}. --- ¡Ah, pero es todavía más divertido ser ebanista que
albañil! Tío, ¿me comprarás un banco, una sierra y un cepillo? Voy a
hacer una caseta para Medor.

%%K t05-01-34 p
\textsc{El Tío}. --- Sí, hijo mío.

%%K t05-01-35 p
\textsc{Juan}. --- ¡Qué contento estoy! ¿Y ese que está junto a la puerta
de entrada, quién es?

%%K t05-tit-6 sub
\VUsaut{3.0mm}
\VUpk{10.2pt}{40}{La pintura.}
\VUsaut{3.0mm}

%%K t05-01-36 p
\textsc{El Sr. Blanco}. --- Es un \VUgras{pintor} \textsuperscript{(102)}. Está en su
escalera con su bote de \VUgras{pintura} \textsuperscript{(103)} y su
\VUgras{brocha} \textsuperscript{(104)}. Pinta la madera de la puerta de entrada
para que se conserve más tiempo.

%%K t05-01-37 p
También es él quien pega el papel en las habitaciones.

%%K t05-01-38 p
El \VUgras{tapicero} \textsuperscript{(107)} que está en una \VUgras{escalera}
\textsuperscript{(106)}, en el \VUgras{balcón} \textsuperscript{(28)}, coloca un
\VUgras{estor} \textsuperscript{(29)}.

%%K t05-01-39 p
El obrero que está junto a la ventana grande es el \VUgras{vidriero}
\textsuperscript{(96)}. Fija los \VUgras{cristales} \textsuperscript{(18)} con
\VUgras{masilla} \textsuperscript{(73)}. Ese otro obrero, arrodillado junto a la
puerta de la bodega, es un \VUgras{cerrajero} \textsuperscript{(97)}. Está poniendo
una \VUgras{cerradura} \textsuperscript{(6)}. Su \VUgras{lima} \textsuperscript{(98)} está a
su lado.

%%K t05-01-40 p
\textsc{Juan}. --- ¿Es también cerrajero ese que golpea con su
\VUgras{martillo} \textsuperscript{(84)} sobre un hierro?

%%K t05-01-41 p
\textsc{El Sr. Blanco}. --- Más bien es un \VUgras{herrero} \textsuperscript{(125)}.
Forja \VUgras{hierros} \textsuperscript{(67)} para el \VUgras{electricista}
\textsuperscript{(124)}, que está en el tejado de la \VUgras{cochera}
\textsuperscript{(51)}. Tiene una \VUgras{fragua} \textsuperscript{(126)}, un
\VUgras{yunque} \textsuperscript{(127)} y unas \VUgras{tenazas} \textsuperscript{(128)}.

%%K t05-01-42 p
El electricista coloca el \VUgras{hilo de cobre} \textsuperscript{(44)} del
teléfono.

%%K t05-tit-7 sub
\VUpk{10.2pt}{40}{La jardinería.}
\VUsaut{3.0mm}

%%K t05-01-43 p
\textsc{El Tío}. --- Al fondo del \VUgras{patio} \textsuperscript{(45)}, junto a la
\VUgras{verja} \textsuperscript{(47)} y al \VUgras{portón} \textsuperscript{(48)}, ves al
\VUgras{jardinero} \textsuperscript{(129)}. Está preparando el pequeño
\VUgras{jardín} \textsuperscript{(49)}. Ha cavado la tierra con su \VUgras{laya}
\textsuperscript{(131)}, siembra semillas y rastrilla con el \VUgras{rastrillo}
\textsuperscript{(130)}. Después, con su \VUgras{regadera} \textsuperscript{(132)}, regará
los \VUgras{arriates} \textsuperscript{(150)} y las plantas crecerán.

%%K t05-01-44 p
El \VUgras{mozo de cuadra} \textsuperscript{(133)}, que acaba de cuidar al caballo y
de darle de comer, limpia el \VUgras{coche} \textsuperscript{(52)} con una esponja,
una \VUgras{bomba de mano} \textsuperscript{(134)} y un cubo de agua. Luego lo
meterá en la cochera y cerrará la puerta con el grueso \VUgras{cerrojo}
\textsuperscript{(53)}.

%%K t05-01-45 p
Ya estás contento, supongo; has tenido bastantes explicaciones.

%%K t05-01-46 p
\textsc{Juan}. --- Sí, tío, pero me gustaría ver el interior de la casa.

%%K t05-01-47 p
\textsc{El Tío}. --- Eso será mañana. El señor Rojo te explicará el plano y
te dirá el nombre de cada habitación.

%%K t05-tit-8 sub
\VUsaut{3.0mm}
\VUpk{10.2pt}{40}{La distribución interior de la casa.}
\VUsaut{2.4mm}

%%K t05-apar-2 apar
{\centering\textit{(Véase el plano.)}\par}
\VUsaut{2.4mm}

%%K t05-01-48 p
\textsc{El Arquitecto}. --- Ven, Juan, voy a hacerte visitar las distintas
partes de la casa.

%%K t05-01-49 p
Entremos en la \VUgras{planta baja} \textsuperscript{(7)}. He aquí el botón del
\VUgras{timbre} \textsuperscript{(11)}. Subimos los tres \VUgras{escalones}
\textsuperscript{(14)} de la \VUgras{escalinata} \textsuperscript{(9)}, pasamos el
\VUgras{umbral} \textsuperscript{(10)} de la \VUgras{puerta de entrada}
\textsuperscript{(8)}, y henos aquí al pie de la \VUgras{escalera} \textsuperscript{(13)}.

%%K t05-01-50 p
\textsc{Juan}. --- Esto es el pasillo.

%%K t05-01-51 p
\textsc{El Arquitecto}. --- No, esto es el \VUgras{vestíbulo} \textsuperscript{(12)}.
El \VUgras{pasillo} \textsuperscript{(a)} está al fondo. A la derecha, he aquí el
\VUgras{salón} \textsuperscript{(b)} y el \VUgras{comedor} \textsuperscript{(c)}; frente al
comedor están la \VUgras{cocina} \textsuperscript{(d)} y la \VUgras{antecocina}
\textsuperscript{(e)}; luego, delante, el \VUgras{despacho} \textsuperscript{(f)}. La
\VUgras{galería} \textsuperscript{(23)}, con su \VUgras{puerta acristalada}
\textsuperscript{(25)}, está junto al salón.

%%K t05-01-52 p
Ahora subiremos la escalera. Agárrate al \VUgras{pasamanos} \textsuperscript{(15)} y
cuenta los escalones. Son catorce. He aquí el \VUgras{rellano}
\textsuperscript{(m)}: estamos en el primer \VUgras{piso} \textsuperscript{(27)}.

%%K t05-tit-9 sub
\VUsaut{3.0mm}
\VUpk{10.2pt}{40}{El primer piso.}
\VUsaut{3.0mm}

%%K t05-01-53 p
Frente a la escalera están el \VUgras{retrete} \textsuperscript{(20)} y el
\VUgras{tocador} \textsuperscript{(j)}. A la derecha, un hermoso \VUgras{dormitorio}
\textsuperscript{(g)} con dos ventanas a la \VUgras{fachada} \textsuperscript{(1)} de la
casa. A la izquierda están el \VUgras{cuarto de baño} \textsuperscript{(1)} y el
\VUgras{cuarto de invitados} \textsuperscript{(i)} con una gran \VUgras{alcoba}
\textsuperscript{(h)}. Junto al dormitorio está el \VUgras{cuarto de los niños}
\textsuperscript{(k)}. Da a una amplia \VUgras{terraza} \textsuperscript{(21)}. Debajo de
esa terraza hay otro retrete \textsuperscript{(20)} y, contra la pared, he aquí la
\VUgras{campana} \textsuperscript{(22)} que llama a cenar.

%%K t05-01-54 p
\textsc{Juan}. --- Subamos al \VUgras{segundo piso}.

%%K t05-01-55 p
\textsc{El Arquitecto}. --- No hay segundo piso. Bajo el tejado solo hay
\VUgras{buhardillas} \textsuperscript{(33)} y el desván \textsuperscript{(34)}. Hemos
terminado la visita, podemos bajar.

%%K t05-01-56 p
\textsc{Juan}. --- Gracias, señor, ha sido muy interesante. Voy a contarle
a mi padre todo lo que he visto.
\VUsaut{4.0mm}
\VUfilet{22mm}
